\section{问题分析与总体技术路线}

本题的四个问题具有明显的递进关系。问题一给出第 5 年末有效需求，是问题二容量配置的输入；问题二确定服务站空间布局和规模，是问题三定价与补贴核算的基础；问题四通过改变关键参数重新求解前三个环节，用于检验方案在不同情景下的稳定性。若直接将四个问题割裂处理，容易出现需求预测、选址容量和定价补贴口径不一致的问题。因此，本文采用统一的数据流和决策流组织模型。

与一般设施选址问题相比，本题的难点有三点。第一，养老服务需求并非静态需求，而是由人口老龄化和失能转移共同驱动。第二，老人满意度并不只由距离决定，还受到响应和价格影响，因此需要将空间可达性与服务体验统一建模。第三，养老服务具有公益属性，不能只追求利润最大化，也不能只追求满意度而忽视机构保本微利。因此，本文将模型目标设置为覆盖优先、容量可行、满意度优化和利润率约束下的综合决策。

为了使模型结果更具解释力，本文在总体方法上设计四个可复核模块。

\begin{enumerate}
    \item \textbf{潜在需求与有效需求分离。} 先计算理论服务需求，再根据消费能力约束得到有效需求，避免后续容量配置按不可支付需求进行虚高建设。
    \item \textbf{感知选择与容量承接分离。} 老人在满意度层面对应唯一主服务站，而服务体系在容量层面允许对实际服务人次进行承接分配，从而同时保留题目中的满意度选择规则和服务能力约束。
    \item \textbf{满意度损失分解。} 将满意度分解为距离、响应和价格三类，使问题二解释空间布局带来的满意度差异，问题三解释价格政策带来的满意度变化。
    \item \textbf{同模型情景重求解。} 问题四不采用口头敏感性讨论，而是重新运行需求预测、选址配置和定价补贴模型，比较站点、价格、补贴、利润率和满意度的联合变化。
\end{enumerate}
